\documentclass[11pt]{article}
\usepackage[margin=1in]{geometry}
\usepackage{amsmath,amssymb,amsthm,hyperref,enumitem}
\hypersetup{colorlinks=true,urlcolor=blue,linkcolor=blue}
\newtheorem{theorem}{Theorem}
\newtheorem{lemma}[theorem]{Lemma}
\newcommand{\pos}[1]{[#1]_+}
\newcommand{\Z}{\mathbb Z}
\title{Positivity of mirror Markov polynomials}
\author{Principia Math In Progress}
\date{22 September 2026}
\begin{document}
\maketitle
\noindent\textbf{MathDB 373161. Task 1: a proof of Conjecture 6.2 of \cite{mirror}.}
For the mirror deformation defined by the initial segment and signed mutations
of \cite[Definitions 4.1--4.2]{mirror}, we prove the following, including the
absence of internal zero coefficients.
\begin{theorem}\label{main}
Every mirror Markov polynomial $m(q)$ has a strictly positive integer coefficient
at each exponent $0,1,\ldots,\deg m$.
\end{theorem}
The input from surface cluster algebras is the established matching expansion
of Musiker--Schiffler--Williams (MSW), together with coefficient specialization.
We give explicitly the cluster realization needed for this application. The
positivity theorem for the \emph{squared} deformation in \cite{mirror} alone
does not imply the result for its factors.

\section*{1. A paired cluster realization}
Use six mutable indices $(i,s)$, where $i\in\{1,2,3\}$ and
$s\in\{+1,-1\}$. Let $B$ have zero diagonal $2\times2$ blocks and all-one
blocks in positions $(1,2),(2,3),(3,1)$, with the opposite blocks their
negatives. Adjoin one frozen variable $q$ and the bottom row
\[
 c_{i,s}=-s. \tag{1}
\]
Initially specialize all six mutable variables to $1$. Ordinary geometric-type
cluster mutation at an index $v$ uses
\[
 x_vx_v'=q^{\pos{c_v}}\prod_{B_{uv}>0}x_u^{B_{uv}}
             +q^{\pos{-c_v}}\prod_{B_{uv}<0}x_u^{-B_{uv}},\qquad
 [a]_+=\max(a,0). \tag{2}
\]
The extended matrix mutates by the usual matrix-mutation formula, as recalled
in \cite[Section 2]{msw}. Define a \emph{paired step at $i$} to mutate both
$(i,+1),(i,-1)$ and then interchange their labels. These two mutations commute.
A direct block calculation shows that each paired step sends $B$ to $-B$:
the two intermediate paths add two arrows against each arrow between the
other two pairs, reversing it after cancellation.

Here is the complete normalization calculation. At any paired seed let $j$
be the other pair with $B_{(j,t),(i,s)}=1$, and let $k$ be the remaining pair.
There are integers $m_i$ and half-integers $u_{i,s}$ such that
\[
 c_{i,s}=m_i-s,\quad z_{i,s}=q^{u_{i,s}}x_{i,s},\quad
 D_i=u_{i,+1}+u_{i,-1},\quad m_i=D_j-D_k. \tag{3}
\]
Initially $m_i=u_{i,s}=D_i=0$. The last equation is
imposed at each pair with its own positive-column neighbor. Put
\[
 a=\pos{m_i+1}+\pos{m_i-1},\qquad
 b=\pos{-m_i-1}+\pos{-m_i+1};\qquad a-b=2m_i.
\]
Matrix mutation and the label interchange give
\[
 m_i'=-m_i,\qquad m_j'=m_j-b,\qquad m_k'=m_k+a. \tag{4}
\]
Before interchanging labels define
\[
 \widetilde u_{i,s}=D_j-\pos{m_i-s}-\frac{s}{2}-u_{i,s},
 \qquad u_{i,s}'=\widetilde u_{i,-s}, \tag{5}
\]
and keep the other $u$'s unchanged. Then $D_i'=2D_j-a-D_i$.
Since the block orientation reverses, the three required new identities in
(3) follow respectively from
\[
 D_k-D_j=-m_i,\quad
 D_i'-D_k=m_j-b,\quad D_j-D_i'=m_k+a.
\]
This proves (3)--(5) inductively for \emph{every} sequence of paired steps.
No restriction on the size or sign of $m_i$ has been used.

Write $A_i=z_{i,+1}z_{i,-1}$. Substituting (3) and (5) into (2), and using
$\pos{v}-\pos{-v}=v$, gives precisely
\[
 z_{i,+1}'=\frac{q^{1/2}A_j+q^{-1/2}A_k}{z_{i,-1}},\qquad
 z_{i,-1}'=\frac{q^{-1/2}A_j+q^{1/2}A_k}{z_{i,+1}}. \tag{6}
\]
Initially $z_{i,+1}=z_{i,-1}=1$. Consequently
$z_{i,-1}(q)=z_{i,+1}(q^{-1})$ throughout. Multiplying the two equations in
(6) gives
\[
 A_iA_i'=A_j^2+(q+q^{-1})A_jA_k+A_k^2. \tag{7}
\]
Equation (6), with the cyclic order reversed at every step, is exactly the
geometric mirror mutation of \cite[Equation (4), Proposition 4.8]{mirror}.
That proposition identifies it, up to powers of $q^{1/2}$, with the signed
polynomial mutation on every branch. For example, the first normalized
polynomials are $1+q$, $1+2q+2q^2$, $2+4q+5q^2+2q^3$, and
$3+8q+10q^2+6q^3+2q^4$, with their reciprocal choices when the initial
orientation is reversed. Thus every polynomial in the question is a monomial
multiple of one of the specialized cluster variables in (2).

\section*{2. The surface and the gap-free matching expansion}
Realize the mutable matrix as the signed adjacency matrix of the tetrahedral
triangulation of the four-punctured sphere. The opposite edge pairs are
\[
 (12,34),\qquad(13,24),\qquad(14,23).
\]
Choose the orientation giving the block matrix above. Each triangular face
contributes one cyclic arrow between the three pairs. Flipping an opposite
pair again gives a tetrahedral triangulation: after the first flip there are
two different arcs between the opposite endpoints, and flipping the original
opposite edge replaces it with an arc between the first endpoints. All four
triangles still have three distinct vertices. In particular, all arcs reached
by paired steps are ordinary arcs, all tagged plain. No notched-arc matching
formula is needed.

Use principal coefficients $y_{i,s}$ temporarily. The MSW expansion
\cite[Theorem 4.9]{msw}, evaluated at all initial mutable variables equal to
$1$, is
\[
 F_\gamma(y)=\sum_{P\in\mathcal M(G_\gamma)}
                   \prod_{T\in H(P)} y_{\ell(T)}. \tag{8}
\]
Here $G_\gamma$ is the ordinary snake graph of the arc, $\mathcal M$ its perfect
matchings, and $H(P)$ the tiles enclosed by the symmetric difference of $P$
and the fixed minimal boundary matching. The label $\ell(T)$ is the initial
arc crossed at that tile. The initial tetrahedral triangulation has no
self-folded triangles, so there are no substitutions in these tile labels.
For an initial arc the expansion is simply $1$.

The coefficient specialization formula
\cite[Theorem 2.13]{msw} identifies our frozen row (1) with
$y_{i,s}=q^{-s}$. Its tropical denominator is a single monomial. It follows
that every specialized $x_\gamma$ in (2) is a monomial multiple of
\[
 \sum_{P\in\mathcal M(G_\gamma)}q^{h(P)},\qquad
 h(P)=\sum_{T\in H(P)}(-s(\ell(T))). \tag{9}
\]
Each tile therefore has height weight $+1$ or $-1$.

\begin{lemma}\label{gaps}
If every tile of an ordinary snake graph has weight $+1$ or $-1$, the
matching sum in (9) has a positive integer coefficient at every integer
between its smallest and largest exponents.
\end{lemma}
\begin{proof}
The perfect matchings of a snake graph are connected by square-tile flips.
For completeness, take an end tile. It adds two vertices to the preceding
snake graph. A matching either contains the edge joining these two new
vertices, or contains both edges joining them to the old vertices. In the
second case, flip that tile, inserting the joining edge and the shared edge.
Now remove the joining edge and the two new vertices. What remains is a
perfect matching of the shorter snake. Induction connects all such matchings,
and their flips lift while keeping the end joining edge fixed. The one-tile
case is immediate.

Under one tile flip the height sets relative to the minimal matching differ
by adding or deleting that one tile. This also follows directly by taking
symmetric differences with its boundary in the definition of $H(P)$.
Consequently $h(P)$ changes by $1$ or $-1$. A path of flips from a matching
of minimum height to one of maximum height visits every intermediate integer.
Every summand has coefficient $1$, which proves the claim.
\end{proof}

\section*{3. Completion of Task 1}
By (3), (6), (9), and Lemma \ref{gaps}, every mirror polynomial has positive
integer coefficients throughout its support interval. The source already
proves that these representatives are polynomials
\cite[Theorem 4.4]{mirror}. We check that their constant terms do not vanish,
to obtain exactly the interval in Theorem \ref{main}.

After the initial segment, order a mirror triple $(a,b,c)$ by its classical
values. Inductively its largest entry $c$ has strictly largest degree and all
constant terms are nonzero. This starts at $(1,1+q,1+2q+2q^2)$.
The corresponding squared Laurent polynomials have the same degrees.
If those degrees are $r\le s<t$, (7) sends a forward mutation to degree
$2t-r$ or $2t-s$, strictly larger than $t$; all numerator coefficients are
nonnegative, so the leading terms cannot cancel. This is also the degree
argument underlying \cite[Proposition 2.3]{mirror}. A new polynomial is,
in the notation of the source,
\[
 c'=q^{\deg c}\frac{q^\varepsilon H+C}{\widetilde d},
 \qquad \varepsilon\in\{-1,1\},\quad
 H=h(q)h(q^{-1}),\quad C=c(q)c(q^{-1}), \tag{10}
\]
where $h$ is the retained smaller entry and $d$ the replaced entry.
Because $\deg c>\deg h$, the lowest exponent of
$q^{\deg c+\varepsilon}H$ is nonnegative. The other numerator term has a
strictly positive constant term, and $\widetilde d$ has a strictly positive
constant term. Thus the polynomial quotient in (10) has a strictly positive
constant term. The initial entries $1,1,1+q$ and $1,1+q,1+2q+2q^2$ satisfy
this explicitly. This closes the induction. The support interval consequently
starts at zero and ends at the degree, proving Theorem \ref{main}.

\paragraph{Attribution and scope.}
The squared deformation, signed recursion, polynomiality, and equivalence
with (6) are established in \cite{mirror}. The surface expansion and
coefficient-specialization theorems are established in \cite{msw}; they are
not new claims here. The explicit frozen row, the all-seed normalization
(3)--(6), and its application to all mirror branches supply the link needed
for the conjecture. Neither the mirror uniqueness conjecture nor the
separate branch-constant conjecture is asserted here.

\paragraph{Verification.}
The accompanying exact verifier checks the tetrahedral matrix, 93 paired
mutation paths of lengths at most five, the normalization and reciprocity at
every step, the printed initial polynomials, and the tile-flip and gap claims
on every north/east snake with at most six tiles and all $\pm1$ tile weights.
These checks support the proof; the finite range is not its justification.
Review is by the authoring agent only, without independent peer review or
proof-assistant verification.

\section*{Proposed follow-up tasks}
\begin{enumerate}[label=Task \arabic*:,start=2,leftmargin=*]
\item \textbf{PROPOSED.} Determine a direct construction of the labeled snake
graph of $m_{a/b}(q)$ from the continued fraction of $a/b$, including its
minimum-height normalization.
\item \textbf{PROPOSED.} Prove or disprove unimodality of the full coefficient
sequence of every mirror Markov polynomial.
\item \textbf{PROPOSED.} Derive closed formulas for the constant and leading
coefficients from the signed left/right mutation word.
\item \textbf{PROPOSED.} Classify integer frozen rows on the tetrahedral quiver
for which every paired mutation variable has a gap-free polynomial
specialization after a monomial normalization.
\end{enumerate}

\begin{thebibliography}{9}
\bibitem{mirror} L. Bittmann, P. Jouteur, E. Kantarc\i{} O\u{g}uz,
M. Molander, and E. Y\i{}ld\i{}r\i{}m,
\emph{A mirror deformation of Markov numbers},
\href{https://arxiv.org/abs/2602.14802}{arXiv:2602.14802v1} (2026).
\bibitem{msw} G. Musiker, R. Schiffler, and L. Williams,
\emph{Positivity for cluster algebras from surfaces},
Adv. Math. \textbf{227} (2011), 2241--2308,
\href{https://doi.org/10.1016/j.aim.2011.04.018}{doi:10.1016/j.aim.2011.04.018}.
Theorem numbering here follows the
\href{https://people.math.harvard.edu/~williams/papers/MSW-July24v5.pdf}{author-hosted version}.
\end{thebibliography}
\end{document}
